% Experimental setup
\section{Experimental Setup}
\label{sec:evaluation}

We evaluate each erasure method by training all probes from scratch on the intervened representations, so that success cannot be attributed to leakage from the erasure scorer itself.

\subsection{Evaluation Protocol}

\paragraph{Fresh-probe evaluation.}
The precision metric is intentionally strict: after every erasure method finishes, we train a \emph{new} probe on the intervened representations. If a freshly trained probe can still recover the concept, then the concept was not actually removed in the sense relevant for downstream use.

\paragraph{Shared validation protocol.}
Most gradient-based methods share a common erasure strength parameter $\lambda$. For \AmbGCE{} and \TanGCE{}, $\lambda$ scales the ambient or tangent-projected update. Because \TanGCE{}'s tangent projection attenuates the effective step size by $\|\mathbf{u}_i^{\mathrm{tan}}\|^2 \leq 1$, a fixed $\lambda$ across methods would systematically under-erase for the tangent-constrained update. Rather than manually choosing different values, we adopt validation-based hyperparameter selection: all $\lambda$-based methods search the same grid and select the $\lambda^*$ that minimizes validation nonlinear probe accuracy. The main comparison (Table~\ref{tab:full_comparison}) and vision benchmarks use $\lambda \in \{0.5, 1.0, 2.0, 4.0, 8.0\}$; the multi-model generalization experiments use the extended grid $\lambda \in \{0.25, 0.5, 1.0, 2.0, 4.0, 8.0, 16.0, 32.0\}$, which we found necessary for larger models and deeper layers where optimal $\lambda^*$ can exceed $8.0$. IGBP has no explicit $\lambda$ parameter---its step size is determined by each sample's distance to the decision boundary---so it runs until convergence.

\paragraph{Training protocol.}
All MLP probes---the shared \emph{concept scorer} used as the erasure direction source, the \emph{validation probe} for $\lambda$ HPO, and the \emph{evaluation probe} retrained on intervened representations---share the same architecture: two-layer MLP with hidden dimension 256, AdamW (lr~$=10^{-3}$, weight decay $10^{-4}$), batch size 512. For language model experiments, all three probes use 1500 gradient steps. For CelebA, the concept scorer uses 1000 steps and the HPO and evaluation probes use 1200 steps; the shorter budget suffices for the lower-dimensional CLIP embeddings. Using identical probe budgets for HPO and evaluation within each modality ensures that the lambda selected on validation is directly comparable to the final test metric, with no ranking distortion from mismatched model capacity. The tangent estimator uses $k{=}50$ nearest neighbors for local PCA; the retained tangent rank per sample is determined by a retained-variance threshold applied to the local eigenspectrum. IGBP runs for up to 35 rounds; each round trains its probe for 10 epochs with AdamW lr~$=2{\times}10^{-4}$, batch size 256, and stops early when validation accuracy falls to majority-class level. Linear baselines (LEACE, INLP, Linear Projection) use $\ell_2$-regularized logistic regression with balanced class weights. The control-task predictor is also an $\ell_2$ logistic regression trained on clean representations and evaluated on the intervened test set without retraining.

\paragraph{Domain-specific controls.}
For sycophancy we use answer matching as the control concept; for Bias in Bios we use profession; for safety we use helpfulness; and for CelebA we use selected low-correlation facial attributes (5 per concept). This lets the evaluation distinguish targeted concept removal from broad representational damage.

\paragraph{Vision benchmarks.}
For cross-modal evaluation we use CelebA~\citep{liu-etal-2015-celeba}, which provides 202{,}599 celebrity images with 40 binary facial attributes; we use the official train/val/test split. We extract pooled CLIP ViT-B/32 image embeddings (768-dim) and apply all erasure methods to the resulting representation matrix.

\subsection{Metrics}

\paragraph{Erasure precision.}
Let $h^{\mathrm{NL}}$ be a nonlinear probe and $h^{\mathrm{L}}$ a linear probe trained on $\tilde{\mathbf{x}}_i$ to predict the target concept $y_i$. We report their test accuracies on the intervened representations:
\begin{equation}
\label{eq:precision_metric}
\mathrm{Prec}_{\mathrm{erase}}^{\mathrm{NL}} = \mathrm{Acc}(h^{\mathrm{NL}}(\tilde{\mathbf{X}}), y),
\qquad
\mathrm{Prec}_{\mathrm{erase}}^{\mathrm{L}} = \mathrm{Acc}(h^{\mathrm{L}}(\tilde{\mathbf{X}}), y).
\end{equation}
Lower values indicate stronger erasure. Because the target concepts in this paper are not assumed to be linearly encoded, the retrained nonlinear probe is the primary precision metric and the linear probe is a secondary diagnostic. For interpretation, the majority-class predictor provides a trivial lower-bound reference for what it means to erase by collapsing the target signal entirely rather than preserving useful structure.

\paragraph{Surgicality on control concepts.}
Let $q$ be a predictor for an annotated control label $c$ (e.g., profession for gender erasure, helpfulness for safety, or selected low-correlation attributes in CelebA). We report the change in control performance relative to the clean representations:
\begin{equation}
\label{eq:surgicality_metric}
\Delta_{\mathrm{ctrl}} = \mathrm{Acc}(q(\tilde{\mathbf{X}}), c) - \mathrm{Acc}(q(\mathbf{X}), c).
\end{equation}
Smaller $|\Delta_{\mathrm{ctrl}}|$ means the erasure is more surgical. In CelebA, where each target attribute is paired with multiple controls, we report the mean $\Delta_{\mathrm{ctrl}}$ across the selected controls.

Together, these metrics ask whether an erasure method removes the target concept while preserving unrelated control concepts.
